% Generated by hostoffload.analysis.pipeline. Do not edit.
\begin{table}[!hb]
\centering
\small
\begin{fittable}{tab:pod-hardware}
\begin{tabular}{lp{2.4cm}p{2.2cm}rp{1.1cm}llp{1.2cm}p{1.6cm}}
\toprule
session & run & CPU & threads & GPU NUMA & GPU link & driver & PCIe max & bus ids \\
\midrule
1 (pod 1) & E8-main/main & AMD EPYC 9534 64-Core & 224 & 1, 1 & NODE & 595.91.07 & Gen5 x16 & 81:00.0, E1:00.0 \\
2 (pod 2) & E12-long-context/long & AMD EPYC 9535 64-Core & 256 & 0, 1 & SYS & 595.91.07 & Gen5 x16 & 06:00.0, 86:00.0 \\
3 (pod 3) & E13-coding-agent/agent & AMD EPYC 9535 64-Core & 256 & 1, 1 & NODE & 595.91.07 & Gen5 x16 & 86:00.0, F4:00.0 \\
\bottomrule
\end{tabular}
\end{fittable}
\caption{Hardware fingerprint of the three pod sessions (all 2x RTX PRO 6000 Blackwell Server Edition, same datacenter), from one run folder per session: CPU model and threads visible to the container (\texttt{fingerprint.json}), NUMA node of each GPU and the path between GPU0 and GPU1 (\texttt{nvidia-smi topo -m}: NODE within one NUMA node, SYS across the inter-socket link), driver, maximum PCIe link and GPU bus ids (\texttt{nvidia-smi -q}).}
\label{tab:pod-hardware}
\end{table}
